\documentclass[aspectratio=169,10pt]{beamer}
\usetheme[numbering=fraction,progressbar=frametitle]{metropolis}
\usepackage{fontspec}
\setsansfont{Latin Modern Sans}
\usepackage{amsmath,amssymb,mathtools,graphicx,booktabs}
\usepackage{polyglossia}
\setdefaultlanguage{french}
\setbeamertemplate{navigation symbols}{}
\setbeamerfont{frametitle}{size=\large}
\setbeamerfont{normal text}{size=\normalsize}
\setlength{\parskip}{.45em}
\hypersetup{pdfauthor={},pdfsubject={Lecture slides},pdfkeywords={},pdftitle={Suites numériques et modèles discrets}}
\providecommand{\N}{\mathbb{N}}
\providecommand{\Z}{\mathbb{Z}}
\providecommand{\D}{\mathbb{D}}
\providecommand{\Q}{\mathbb{Q}}
\providecommand{\R}{\mathbb{R}}
\providecommand{\C}{\mathbb{C}}
\providecommand{\vect}[1]{\overrightarrow{#1}}
\providecommand{\vi}{\vec{\imath}}
\providecommand{\vj}{\vec{\jmath}}
\providecommand{\vk}{\vec{k}}
\providecommand{\norm}[1]{\left\lVert #1\right\rVert}
\providecommand{\vc}[2]{\begin{pmatrix}#1\\#2\end{pmatrix}}
\providecommand{\deter}[2]{\det\left(#1,#2\right)}
\providecommand{\abs}[1]{\left\lvert #1\right\rvert}
\providecommand{\dd}{\,\mathrm{d}}
\providecommand{\e}{\mathrm{e}}
\providecommand{\Card}{\operatorname{Card}}
\providecommand{\E}{\operatorname{E}}
\providecommand{\V}{\operatorname{V}}
\providecommand{\degre}{^{\circ}}
\providecommand{\un}[1]{\,\mathrm{#1}}
\providecommand{\sci}[2]{#1\times10^{#2}}
\providecommand{\pgcd}{\operatorname{PGCD}}
\providecommand{\ppcm}{\operatorname{PPCM}}
\title{Suites numériques et modèles discrets}
\subtitle{Cours et exemples}
\date{}
\author{}
\begin{document}
\begin{frame}[plain]
\vfill
{\usebeamerfont{title}\inserttitle\par}
\medskip
{\usebeamerfont{subtitle}\insertsubtitle\par}
\medskip
{\usebeamercolor[fg]{progress bar}\rule{\textwidth}{.5pt}}
\vfill
\end{frame}
\begin{frame}{Le parcours du chapitre}
\textbf{Objectifs.} Passer du récit à une suite ; distinguer accroissement et taux ; calculer des termes et des sommes ; justifier une variation ; expérimenter un seuil.

\textbf{Prérequis.} Calcul littéral, puissances entières, pourcentages, fonction affine, boucle Python.

\textbf{Repère annuel.} 8-17 septembre : 8 h, puis reprises
\end{frame}
\begin{frame}{Définir et représenter une suite}
Une suite réelle est une fonction dont la variable est un entier : à chaque rang \(n\) autorisé correspond un réel \(u_n\). La notation \((u_n)\) désigne la suite, \(u_n\) un terme. Une formule explicite donne directement \(u_n\) en fonction de \(n\). Une définition par récurrence fournit un premier terme et une règle de passage. La représentation est un ensemble de points \((n;u_n)\) : on ne relie pas nécessairement ces points.
\end{frame}
\begin{frame}{Définir et représenter une suite}
\begin{columns}[T,onlytextwidth]
\begin{column}{.57\textwidth}
Pour \(u_n=n^2+1\), \(u_0=1\), \(u_3=10\). Pour \(v_0=1\) et \(v_{n+1}=2v_n+1\), \(v_1=3\), \(v_2=7\), \(v_3=15\). On ne peut pas remplacer \(v_3\) par \(2\times3+1\) : la règle porte sur le terme précédent.
\end{column}
\begin{column}{.4\textwidth}
\centering
\includegraphics[width=\linewidth,height=4.5cm,keepaspectratio]{figures/premiere-01-1.pdf}
\end{column}
\end{columns}
\end{frame}
\begin{frame}{Définir et représenter une suite}
Le rang n’est pas la valeur. Une liste finie de termes ne détermine pas une règle unique.
\end{frame}
\begin{frame}{Définir et représenter une suite}
Calculer les quatre premiers termes de \(w_0=2\), \(w_{n+1}=w_n+2n+1\) ; proposer une formule explicite.
\end{frame}
\begin{frame}{Suites arithmétiques}
Une suite est arithmétique de raison \(r\) si \(u_{n+1}=u_n+r\) pour tout rang autorisé. Alors \(u_n=u_p+(n-p)r\) pour \(n\geq p\). Si \(r>0\), elle est strictement croissante ; si \(r<0\), strictement décroissante ; si \(r=0\), constante. Les points sont alignés sur une droite.
\end{frame}
\begin{frame}{Suites arithmétiques}
Une réserve contient 140 L et reçoit 12 L par minute. Après \(n\) minutes, \(u_n=140+12n\). À 5 min : 200 L. Si l’on commence l’indexation à 1 avec 140 L, on écrit \(u_n=140+12(n-1)\).
\end{frame}
\begin{frame}{Suites arithmétiques}
Ajouter un pourcentage constant ne donne généralement pas une suite arithmétique.
\end{frame}
\begin{frame}{Suites arithmétiques}
On sait \(u_3=17\) et \(u_8=37\) pour une suite arithmétique. Déterminer \(r\) et \(u_0\).
\end{frame}
\begin{frame}{Suites géométriques}
Une suite est géométrique de raison \(q\) si \(u_{n+1}=qu_n\). Pour \(n\geq p\), \(u_n=u_pq^{n-p}\), avec la convention qu’un produit vide vaut 1. Un taux constant \(t\) donne \(q=1+t\). La variation se décide avec \(u_{n+1}-u_n=(q-1)u_n\). Pour \(q<0\), les signes peuvent alterner.
\end{frame}
\begin{frame}{Suites géométriques}
Une quantité de 160 mg perd 25 \% à chaque étape : \(u_n=160(\dfrac{3}{4})^n\). Les termes sont positifs et \(u_{n+1}-u_n=-\dfrac{u_n}{4}<0\). On obtient \(u_2=90\) mg, sans confondre baisse de 25 \% et retrait de 25 mg.
\end{frame}
\begin{frame}{Suites géométriques}
\(q>1\) ne suffit pas à conclure que toute suite géométrique est croissante : le signe du premier terme compte.
\end{frame}
\begin{frame}{Suites géométriques}
Étudier la variation de \(u_n=-2\times3^n\), puis de \(v_n=2(-1)^n\).
\end{frame}
\begin{frame}{Sommes, indices et démonstrations}
Pour \(n\geq1\), \(1+2+\cdots +n=\dfrac{n(n+1)}{2}\). Pour \(q\ne1\) et \(n\geq0\), \(1+q+\cdots +q^n=\dfrac{1-q^{n+1}}{1-q}\). Pour \(q=1\), la somme vaut \(n+1\). Ainsi, de \(u_p\) à \(u_n\), une somme arithmétique vaut \(\dfrac{(n-p+1)(u_p+u_n)}{2}\) ; une somme géométrique vaut \(\dfrac{u_p(1-q^{n-p+1})}{1-q}\) si \(q\ne1\).
\end{frame}
\begin{frame}{Sommes, indices et démonstrations}
Pour \(u_n=3+2n\), \(u_2+\cdots +u_6=\dfrac{5(7+15)}{2}=55\). Pour \(v_n=2\times3^n\), \(v_0+\cdots +v_3=\dfrac{2(1-3^4)}{1-3}=80\).
\end{frame}
\begin{frame}{Sommes, indices et démonstrations}
Compter les termes avant d’utiliser la formule : de 0 à n, il y en a n+1.
\end{frame}
\begin{frame}{Sommes, indices et démonstrations}
Calculer \(\sum_{k=2}^{5}2^k\) sans symbole de sommation dans la réponse.
\end{frame}
\begin{frame}{Conjecturer une limite et programmer un seuil}
Une limite est ici abordée intuitivement. Les termes \(\dfrac{1}{n+1}\) se rapprochent de 0 ; \(n^2\) grandit sans borne ; \((-1)^n\) ne se stabilise pas. Une expérience numérique suggère une conclusion mais ne prouve pas un comportement pour tous les rangs. Un algorithme de seuil doit préciser le sens de l’inégalité et prévoir un arrêt.
\end{frame}
\begin{frame}{Conjecturer une limite et programmer un seuil}
Pour \(u_n=80(\dfrac{1}{2})^n\), le premier rang tel que \(u_n<6\) est 4 : \(u_3=10\geq6\) et \(u_4=5<6\). Le programme commence avec \(n=0,u=80\) et répète la division par 2 tant que \(u\geq6\).
\end{frame}
\begin{frame}{Conjecturer une limite et programmer un seuil}
Une boucle sans borne peut ne jamais s’arrêter si le seuil est inaccessible ; \(<\) et \(\leq\) ne sont pas interchangeables.
\end{frame}
\begin{frame}{Conjecturer une limite et programmer un seuil}
\begin{columns}[T,onlytextwidth]
\begin{column}{.57\textwidth}
Pour \(u_n=48(\dfrac{1}{2})^n\), déterminer les premiers rangs où \(u_n\leq6\) et où \(u_n<6\).
\end{column}
\begin{column}{.4\textwidth}
\centering
\includegraphics[width=\linewidth,height=4.5cm,keepaspectratio]{figures/premiere-01-2.pdf}
\end{column}
\end{columns}
\end{frame}
\end{document}
